\clearpage
\VerdierPhysicalPageAnchor{0193}%
\begingroup\small\setlength{\parskip}{0pt}\setlength{\emergencystretch}{1.5em}
\setlength{\abovedisplayskip}{4pt plus 2pt minus 2pt}%
\setlength{\belowdisplayskip}{4pt plus 2pt minus 2pt}%

\VerdierSemanticAnchor{ch3d-2-2-6-proof-cont}
(2.1.1). Le morphisme $p_0$ est un morphisme élémentaire de hauteur $1$.
Le complexe $U_0^\bullet$ est de la forme :

\VerdierSemanticAnchor{ch3d-2-2-6-u0-complex}
\[
\begin{aligned}
U_0^\bullet={}&0\longrightarrow0\longrightarrow\cdots\longrightarrow0
\longrightarrow L_n\longrightarrow L_{n-1}\longrightarrow\cdots\\
&\longrightarrow L_1\longrightarrow Z_0\longrightarrow L
\longrightarrow0\longrightarrow0\longrightarrow\cdots.
\end{aligned}
\]
D'après la propriété $(E_4)$, l'objet $Z_0$ est un élément de $M'$. Par
suite, il existe un épimorphisme $u_0\colon L'_0\longrightarrow Z_0$,
$L'_0\in M$. Posons alors (2.1.1)
$U_1^\bullet=U_0^\bullet(u_0,0)$,
$p_1=p(u_0,0)\colon U_1^\bullet\longrightarrow U_0^\bullet$.

\VerdierSemanticAnchor{ch3d-2-2-6-iterated-elementary-morphisms}
En poursuivant ainsi, on démontre par récurrence, qu'il existe une suite de
morphismes :
\[
p_i\colon U_i^\bullet\longrightarrow U_{i-1}^\bullet,
\qquad 1\leq i\leq n-1,
\]
\[
p_0\colon U_0^\bullet\longrightarrow V^\bullet,
\]
où $p_i$ est un morphisme élémentaire de hauteur $-i+1$ et où
$U_i^\bullet$ est de la forme :

\VerdierSemanticAnchor{ch3d-2-2-6-ui-complex}
\[
\begin{aligned}
U_i^\bullet={}&0\longrightarrow0\longrightarrow\cdots
\longrightarrow L_n\longrightarrow L_{n-1}\longrightarrow\cdots
\longrightarrow L_{i+1}\longrightarrow Z_i\\
&\longrightarrow L'_{i-1}\longrightarrow\cdots\longrightarrow L'_0
\longrightarrow L\longrightarrow0\longrightarrow\cdots ,
\end{aligned}
\]
où $L'_j\in M$, $0\leq j\leq i-1$. On obtient donc, en composant ces
différents morphismes $p_i$, un quasi-isomorphisme épimorphique :

\VerdierSemanticAnchor{ch3d-2-2-6-composite}
\[
p\colon U_{n-1}^\bullet\longrightarrow V^\bullet.
\]
Remarquons alors que dans le complexe :

\VerdierSemanticAnchor{ch3d-2-2-6-unminus1-complex}
\[
\begin{aligned}
U_{n-1}^\bullet={}&0\longrightarrow0\longrightarrow\cdots
\longrightarrow L_n\longrightarrow Z_{n-1}\longrightarrow L'_{n-2}
\longrightarrow\cdots\\
&\longrightarrow L'_0\longrightarrow L\longrightarrow0
\longrightarrow\cdots ,
\end{aligned}
\]
l'objet $Z_{n-1}$ apparaît, d'après la propriété $(E_4)$, comme une
extension de deux éléments de $M$ et est par suite, d'après la propriété
$(E_5)$, un élément de $M$. Posons alors :

\VerdierSemanticAnchor{ch3d-2-2-6-l-complexes}
\[
L^\bullet=0\longrightarrow0\longrightarrow L_n\longrightarrow L_{n-1}
\longrightarrow\cdots\longrightarrow L_0\longrightarrow0
\longrightarrow0\longrightarrow\cdots ,
\]
\[
\begin{aligned}
{L'}^\bullet={}&0\longrightarrow0\longrightarrow L_n
\longrightarrow Z_{n-1}\longrightarrow L'_{n-2}
\longrightarrow\cdots\\
&\longrightarrow L'_0\longrightarrow0\longrightarrow0
\longrightarrow\cdots ,
\end{aligned}
\]
et notons $m\colon {L'}^\bullet\longrightarrow L^\bullet$ le morphisme
induit par $p$. Le seul objet de cohomologie non nul de $L^\bullet$ (resp.
${L'}^\bullet$) est l'objet $Y$ en dimension $0$ (resp. $L$ en dimension
$0$) et le morphisme $m$ induit sur ces objets le morphisme $u$.
Considérons alors le cône $c(m)$ du morphisme $m$ :

\VerdierSemanticAnchor{ch3d-2-2-6-cone}
\[
\begin{aligned}
c(m)={}&\cdots\longrightarrow0\longrightarrow L_n
\longrightarrow L_n\oplus Z_{n-1}
\longrightarrow L_{n-1}\oplus L'_{n-2}\longrightarrow\cdots\\
&\longrightarrow L_1\oplus L'_0\longrightarrow L_0
\longrightarrow0\longrightarrow\cdots.
\end{aligned}
\]
\endgroup
